% GENERATED FILE. Edit canonical lessons, then run npm run generate:books.
% canonical-source-hash: fnv1a64:f083bf9f2b2a3681
% canonical-lessons: GE-C310-karneval, GE-C310-fasching, GE-C310-kostuem, GE-C310-verkleiden, GE-C310-brauch

\chapter{Karneval}
\label{ch:ge-karneval}
\hlchaptermodality{310}

\begin{chapteropening}
\textbf{By the end of this chapter:} \emph{I can talk about carnival: carnival, Fasching, a costume, to dress up and a custom.}

Complete the last lesson of chapter 310: I can talk about carnival: carnival, Fasching, a costume, to dress up and a custom.
\end{chapteropening}

\section[der Karneval]{der Karneval — the carnival}

\label{lesson:GE-C310-karneval}



\begin{quote}
Before the new one: say the German for the balloon, then the German for the fireworks.
\end{quote}

\subsection*{You'll want to know: der Karneval}
\textbf{der Karneval} — \textquotedblleft{}the carnival\textquotedblright{}.

\begin{quote}
Im Februar feiert die Stadt Karneval.
\end{quote}

\begin{quote}
\emph{In February the town celebrates carnival.}
\end{quote}

Plural: \emph{die Karnevale}. The name differs by region; in the southwest it is \emph{die Fastnacht}.

\begin{cousinweb}
\emph{Karneval} comes from Italian \emph{carnevale}, from Medieval Latin \emph{carnelevale}, built on \emph{caro} \textquotedblleft{}flesh, meat\textquotedblright{} and \emph{levāre} \textquotedblleft{}to lighten, to raise\textquotedblright{}, here \textquotedblleft{}to take away\textquotedblright{} (EtymDB): the farewell to meat before Lent.
\end{cousinweb}

\begin{grammarlens}[title={What you've built}]
One more word for telling, explaining and describing at B1.
\end{grammarlens}

\subsection*{Your turn}
\begin{itemize}
\raggedright
  \item \emph{Say it:} \emph{der Karneval}
  \item \emph{Say it:} \emph{der Karneval}, once more
  \item \emph{Say it:} a sentence of your own with \emph{der Karneval}
  \item \emph{Recall:} say the German for the balloon, then the German for the fireworks, then say \emph{der Karneval} again
\end{itemize}

\begin{tcolorbox}[breakable,colback=teal!4,colframe=teal!35!black,title={Before you move on}]
What does der Karneval mean? (\textquotedblleft{}The carnival\textquotedblright{}.) Say it once more.
\end{tcolorbox}
\section[der Fasching]{der Fasching — the Shrovetide carnival}

\label{lesson:GE-C310-fasching}



\begin{quote}
Before the new one: say the German for the fireworks, then the German for the carnival.
\end{quote}

\subsection*{You'll want to know: der Fasching}
\textbf{der Fasching} — \textquotedblleft{}the Shrovetide carnival\textquotedblright{}.

\begin{quote}
Im Fasching gibt es in der Schule eine Party.
\end{quote}

\begin{quote}
\emph{At carnival time there is a party at school.}
\end{quote}

The word in southern Germany and Austria; it is the season before Lent.

\begin{cousinweb}
\emph{Fasching} is usually traced to Middle High German \emph{vaschanc} \textquotedblleft{}the serving of drink before the fast\textquotedblright{}; EtymDB has the word but lists no parent.
\end{cousinweb}

\begin{grammarlens}[title={What you've built}]
One more word for telling, explaining and describing at B1.
\end{grammarlens}

\subsection*{Your turn}
\begin{itemize}
\raggedright
  \item \emph{Say it:} \emph{der Fasching}
  \item \emph{Say it:} \emph{der Fasching}, once more
  \item \emph{Say it:} a sentence of your own with \emph{der Fasching}
  \item \emph{Recall:} say the German for the fireworks, then the German for the carnival, then say \emph{der Fasching} again
\end{itemize}

\begin{tcolorbox}[breakable,colback=teal!4,colframe=teal!35!black,title={Before you move on}]
What does der Fasching mean? (\textquotedblleft{}The Shrovetide carnival\textquotedblright{}.) Say it once more.
\end{tcolorbox}
\section[das Kostüm]{das Kostüm — the costume; a skirt suit}

\label{lesson:GE-C310-kostuem}



\begin{quote}
Before the new one: say the German for the carnival, then the German for the Shrovetide carnival.
\end{quote}

\subsection*{You'll want to know: das Kostüm}
\textbf{das Kostüm} — \textquotedblleft{}the costume; a skirt suit\textquotedblright{}.

\begin{quote}
Sie nähte ihr Kostüm selbst.
\end{quote}

\begin{quote}
\emph{She sewed her costume herself.}
\end{quote}

Plural: \emph{die Kostüme}. Also a woman's skirt suit.

\begin{cousinweb}
\emph{Kostüm} was borrowed from French \emph{costume}, which came from Italian and goes back to Latin \emph{cōnsuētūdō} \textquotedblleft{}custom\textquotedblright{} (EtymDB): first the dress of a time and place.
\end{cousinweb}

\begin{grammarlens}[title={What you've built}]
One more word for telling, explaining and describing at B1.
\end{grammarlens}

\subsection*{Your turn}
\begin{itemize}
\raggedright
  \item \emph{Say it:} \emph{das Kostüm}
  \item \emph{Say it:} \emph{das Kostüm}, once more
  \item \emph{Say it:} a sentence of your own with \emph{das Kostüm}
  \item \emph{Recall:} say the German for the carnival, then the German for the Shrovetide carnival, then say \emph{das Kostüm} again
\end{itemize}

\begin{tcolorbox}[breakable,colback=teal!4,colframe=teal!35!black,title={Before you move on}]
What does das Kostüm mean? (\textquotedblleft{}The costume; a skirt suit\textquotedblright{}.) Say it once more.
\end{tcolorbox}
\section[verkleiden]{verkleiden — to dress up (as); mostly reflexive: sich verkleiden}

\label{lesson:GE-C310-verkleiden}



\begin{quote}
Before the new one: say the German for the Shrovetide carnival, then the German for the costume.
\end{quote}

\subsection*{You'll want to know: verkleiden}
\textbf{verkleiden} — \textquotedblleft{}to dress up (as); mostly reflexive: sich verkleiden\textquotedblright{}.

\begin{quote}
Mein Bruder verkleidete sich als Löwe.
\end{quote}

\begin{quote}
\emph{My brother dressed up as a lion.}
\end{quote}

\emph{Sich als etwas verkleiden}. \emph{Die Verkleidung} — the disguise.

\begin{cousinweb}
\emph{verkleiden} is built on \emph{Kleid} \textquotedblleft{}dress\textquotedblright{}, an inherited Germanic word with the Dutch cousin \emph{kleed} (EtymDB), with \emph{ver-}.
\end{cousinweb}

\begin{grammarlens}[title={What you've built}]
One more word for telling, explaining and describing at B1.
\end{grammarlens}

\subsection*{Your turn}
\begin{itemize}
\raggedright
  \item \emph{Say it:} \emph{verkleiden}
  \item \emph{Say it:} \emph{verkleiden}, once more
  \item \emph{Say it:} a sentence of your own with \emph{verkleiden}
  \item \emph{Recall:} say the German for the Shrovetide carnival, then the German for the costume, then say \emph{verkleiden} again
\end{itemize}

\begin{tcolorbox}[breakable,colback=teal!4,colframe=teal!35!black,title={Before you move on}]
What does verkleiden mean? (\textquotedblleft{}To dress up (as); mostly reflexive: sich verkleiden\textquotedblright{}.) Say it once more.
\end{tcolorbox}
\section[der Brauch]{der Brauch — the custom}

\label{lesson:GE-C310-brauch}



\begin{quote}
Before the new one: say the German for the costume, then the German for to dress up.
\end{quote}

\subsection*{You'll want to know: der Brauch}
\textbf{der Brauch} — \textquotedblleft{}the custom\textquotedblright{}.

\begin{quote}
Das ist ein alter Brauch in unserem Dorf.
\end{quote}

\begin{quote}
\emph{That is an old custom in our village.}
\end{quote}

Plural: \emph{die Bräuche}.

\begin{cousinweb}
\emph{Brauch} comes from Middle High German \emph{brūch}, Old High German \emph{brūh} \textquotedblleft{}use\textquotedblright{} (EtymDB), related to \emph{brauchen}: what people do by use.
\end{cousinweb}

\begin{grammarlens}[title={What you've built}]
One more word for telling, explaining and describing at B1.
\end{grammarlens}

\subsection*{Your turn}
\begin{itemize}
\raggedright
  \item \emph{Say it:} \emph{der Brauch}
  \item \emph{Say it:} \emph{der Brauch}, once more
  \item \emph{Say it:} a sentence of your own with \emph{der Brauch}
  \item \emph{Recall:} say the German for the costume, then the German for to dress up, then say \emph{der Brauch} again
\end{itemize}

\begin{tcolorbox}[breakable,colback=teal!4,colframe=teal!35!black,title={Before you move on}]
What does der Brauch mean? (\textquotedblleft{}The custom\textquotedblright{}.) Say it once more.
\end{tcolorbox}
